\section{How a saved field becomes a screen image}
\subsection{Native reading and the process boundary}
React and TypeScript own interaction, view state and rendering preparation.
Tauri's Rust layer owns native filesystem access, readers and process operations.
Large numeric payloads cross the native interface as binary arrays with explicit
metadata rather than as enormous JSON lists. The reader validates dimensions,
types, bounds and source identities before the frontend treats the payload as a
field.

A bounded-window request describes an origin, extent and sampling policy.
Returned bytes and source work are separate quantities. A block mean can return
one value while visiting many source cells; a compressed representation can need
codec scratch in addition to its decoded output. Accounting only for the final
array would underestimate the temporary native read footprint.

The native netCDF API is loaded once and retained for the process lifetime,
while individual file handles are still closed after their operations. This
distinguishes the lifetime of a library's executable code and internal cleanup
state from the lifetime of a dataset file. Repeatedly unloading a library that
has registered process or thread cleanup callbacks can leave callbacks referring
to unloaded code. Keeping the library available does not mean retaining every
opened dataset or bypassing the reader's serialization and memory controls.

\subsection{Workers prepare data without owning the interface}
Web workers perform expensive preparation, derived-field work, contour
extraction, streamline integration and other bounded tasks outside the UI thread.
The interface can remain responsive while a worker scans a field or builds a
mesh. Moving work to a worker does not make it free: memory, transfers, queues
and completion synchronization still contribute to latency.

Array buffers can be transferred between owners rather than repeatedly copied.
A transfer changes who may use the buffer; the sender cannot continue treating a
detached buffer as its live backing. Worker dataset registries therefore record
identity and lifetime, not just a pointer or field label. Geometry caches and
scalar caches have to agree on the actual resident dataset used by a task.

Changing a field or frame creates a new request generation. A completed old task
may no longer be useful when it returns. The result is validated against the
current owner before publication, so late work cannot overwrite the newly
selected view. Cancellation also retires obsolete resource claims. A cache may
retain a deliberately reusable result, but that retention must have an owner and
be represented in the relevant accounting.

\subsection{Plotly and VTK have different jobs}
Plotly handles most two-dimensional figures, profiles, sounding diagrams,
history plots, Python numeric previews and research panels. Its canvas and vector
elements manage axes, labels, selection and zoom. VTK handles the primary
volumetric, surface, streamline and trajectory scenes. It consumes prepared
datasets and geometry, constructs actors and mappers, and submits graphics work
through the chosen browser graphics API.

These are application-rendering paths. The solver's Vulkan and Metal compute
backends do not determine Studio's graphics API. A CPU simulation can be viewed
with GPU rendering; a GPU simulation can be viewed on another computer. The
desktop webview, driver, graphics capability and workload determine which
rendering paths are available.

Volume workloads currently use the WebGL GPU path because the maintained
vtk.js WebGPU volume implementation does not provide the same required shading,
opacity-normal, scattering and close-camera behavior. Geometry workloads can
select WebGPU when adapter discovery and the actual VTK device/context
initialization succeed. A preference alone is not proof that the requested API
became active. Device or initialization failure is disclosed and uses the
maintained WebGL path where available. This rendering fallback is separate from
the solver's native-required compute admission contract.

\subsection{Textures, triangles and shaders}
For a surface or streamline, the CPU or worker prepares vertices and connectivity.
The graphics API uploads geometry, transforms vertices through the scene and
camera matrices, rasterizes primitives and shades their covered pixels. Surface
normals and light directions affect brightness. A bright face can be facing a
light rather than carrying a larger field value; the colorbar describes the
scalar mapping, not the lighting multiplier.

For a volume, a shader samples one or more scalar textures along each viewing
ray. Each sample is mapped to a color and opacity. A simplified front-to-back
composition illustrates the ownership of those choices:
\begin{align}
 C_{j+1} &= C_j+(1-A_j)\alpha_j c_j,\\
 A_{j+1} &= A_j+(1-A_j)\alpha_j.
\end{align}
Here $c_j$ is a sample color, $\alpha_j$ its opacity over the current sample
segment, and $C_j,A_j$ the accumulated color and opacity. Both the step length
and transfer function matter. A physically scaled opacity convention can be
expressed as $\alpha=1-\exp(-\sigma\,\Delta s)$ for an effective attenuation
$\sigma$ and path increment $\Delta s$. This explains why changing sampling
distance without correcting opacity can change how thick a cloud appears.
It is an explanation of rendering, not a replacement for the solver's
radiation calculation.

Texture interpolation, scalar validity, lighting, shadow sampling, background
composition and output color are separate shader responsibilities. Shadows help
depth perception but do not represent the atmospheric heating computed by a
radiation scheme. A reflectivity floor is composited with its own color encoding
and physical plane. The final framebuffer holds rendered pixels; the native
field and its scientific units remain separate data.

\subsection{Context loss and scene reconstruction}
Graphics contexts can be lost because of driver resets, resource pressure or
the browser's management of graphics resources. Once a context is lost, its
old textures and framebuffer attachments cannot be assumed usable. Studio's
owned render-window construction installs traversal protection and framebuffer
ownership together before the VTK object is frozen. Independent patches must
not overwrite one another's guards.

The loss path stops invalid traversal; restoration reconstructs the scene through
its owners. A restored canvas must reach a completed render before it is reported
as a current frame or captured. Merely receiving a context-restored event is not
enough. This boundary prevents a stale texture, half-rebuilt layer or previous
frame's legend from masquerading as the selected result.

Recovery cannot guarantee that every driver or device can support the requested
scene. Persistent failure should remain visible and can require a smaller scene,
different graphics preference or application restart. What the implementation
must preserve is honest state: it cannot silently label invalid or incomplete
pixels as a completed scientific view.

\section{Frame identity, playback and coherent presentation}
\subsection{Requested, loaded and displayed are distinct}
Selecting a saved frame starts a request; it does not immediately replace the
pixels. A field can finish reading before its geometry is prepared, and a scene
can finish preparation before the graphics API draws it. Studio tracks these
stages independently. The displayed time should describe the pixels currently
visible, with the requested frame shown separately when it differs.

Frame identity includes both the saved step and its publication generation.
Reopening a reused directory can produce another run with the same step number.
Matching only that number would accept data from the wrong publication. Multi-
field plots require every contributing array to agree on the exact publication,
not just on shape or step ordinal.

The capture predicate is therefore conceptually
\begin{equation}
 \text{capture allowed} =
 (\text{displayed identity}=\text{requested identity})
 \ \land\ (\text{render complete}).
\end{equation}
The implementation additionally guards the active source and operation while
capture proceeds. A late change in field, view, source or frame invalidates work
that no longer describes the current request.

\begin{figure}[htbp]
\centering
\begin{tikzpicture}[
  box/.style={draw=NimbusBlue,rounded corners,fill=NimbusLight,align=center,
    text width=2.5cm,minimum height=1.1cm,font=\small},
  flow/.style={-{Latex[length=2mm]},thick,draw=NimbusBlue}]
\node[box] (request) {Requested\\step and generation};
\node[box,right=0.5cm of request] (load) {Loaded\\matching fields};
\node[box,right=0.5cm of load] (prepare) {Prepared\\coherent layers};
\node[box,right=0.5cm of prepare] (display) {Displayed\\pixels and receipt};
\draw[flow] (request) -- (load);
\draw[flow] (load) -- (prepare);
\draw[flow] (prepare) -- (display);
\node[below=0.6cm of load,align=center,font=\small,text width=5.4cm] (reject)
  {Cancellation, source change or identity mismatch\\retires obsolete work before publication};
\draw[flow,dashed] (load.south) -- (reject.north);
\node[below=0.6cm of display,align=center,font=\small,text width=3.0cm] (capture)
  {Capture requires the exact\\completed identity};
\draw[flow] (display) -- (capture);
\end{tikzpicture}
\caption{A request is not a displayed frame. The last completed scene may remain
visible during preparation; its label and capture state must continue to describe
that scene until the new one commits.}
\end{figure}


\subsection{Multi-layer scenes commit together}
A cloud volume, updraft surface and radar floor can have different preparation
costs. Publishing each as soon as it happens to finish could mix atmospheric
times in one picture. A coordinated scene keeps the selected generation and
layer identities until the required work is complete. Legends, source disclosure
and capture metadata are committed with the scene they describe.

Appearance-only updates have a smaller boundary: they retain the scientific
frame and geometry, redraw the changed material, then refresh the appearance
part of the capture descriptor. Updating a shadow control must not accidentally
claim that a new scientific field was loaded. Conversely, retaining the right
frame is not permission to leave the recorded lighting recipe stale.

\subsection{Playback is an output-viewing schedule}
Playback advances through saved frames, waiting for the participating views'
required completion contracts. The requested rate is a target; achievable
cadence depends on source I/O, decoding, preparation and rendering. A 60 fps
request does not create model states between saved outputs and does not promise
60 newly rendered scientific frames per second on every machine.

Prefetch can prepare future reads while the current frame is displayed. It must
remain subordinate to visible work and resource ownership. Leaving a plot or
deselecting its field retires its unused prefetch requests. A shared request
continues when another visible consumer still needs it. Canceling one pane must
not break another pane's valid use of the same field.

The distinction between startup and steady interaction also matters. Creating a
large first scene can take longer than changing an already prepared view.
Measuring responsiveness before the first completed scientific frame describes
startup, not steady playback. Both are useful measurements, but they should be
reported as different intervals rather than merged into an unsupported speed
claim.
